\documentclass[11pt]{article}

\usepackage[T1]{fontenc}
\usepackage[margin=1in]{geometry}
\usepackage{amsmath, amssymb, amsthm}
\usepackage{enumitem}
\usepackage{hyperref}
\usepackage{booktabs}
\usepackage{listings}
\usepackage{xcolor}

\lstset{
  basicstyle=\ttfamily\small,
  frame=single,
  breaklines=true,
  columns=fullflexible,
  mathescape=true
}

\theoremstyle{definition}
\newtheorem{theorem}{Theorem}[section]
\newtheorem{definition}[theorem]{Definition}
\newtheorem{example}[theorem]{Example}

\title{CS 250 --- Module 5: Informal Proofs and Groups}
\author{}
\date{}

\begin{document}
\maketitle

This module is, in a sense, a breather. We're not introducing new concepts here --- instead we're getting practice with informal proofs after several modules of formal logic.

The relevant sections of Epp give us a lot of time to practice proofs with numbers. They're very well written and even call out some topics I think are important: a games-based approach to quantifiers, constructive vs.\ non-constructive proofs of $\exists$, and the relationship between writing proofs and writing computer programs (framing the connection as analogy rather than mathematical identity).

What I think might be useful as an addition to the text, rather than attempting to recapitulate its content, is to give lots of examples of proofs in a different context than numbers.

But first, let's make sure we have the number-theory toolkit down solid, since that's what Epp 4.1--4.4 is all about.

One thing worth saying up front: the informal proof techniques we're about to practice are not a departure from the formal rules we built up in Modules 3 and 4---they're implementations of them. A direct proof of ``for all $x$, if $P(x)$ then $Q(x)$'' is really $\forall$-intro combined with $\to$-intro: introduce a fresh variable, assume the hypothesis, derive the conclusion. A counterexample to a universal claim is $\exists$-intro applied to the negation $\lnot(\forall x.\; P(x))$. A proof by cases is $\lor$-elimination in disguise. The formal rules tell us \emph{why} these techniques are valid; the informal style is how mathematicians actually write. We'll keep pointing out these connections as we go, so you can see the formal machinery underneath the prose.

\section{Direct proof}

The most basic proof technique is the \emph{direct proof}\index{proof!direct}. The shape of a direct proof for a universal conditional statement --- something of the form ``for all $x$, if $P(x)$ then $Q(x)$'' --- is: you assume $P(x)$ for an \emph{arbitrary} $x$, and then you show that $Q(x)$ follows. That's it. The whole game is figuring out how to get from $P(x)$ to $Q(x)$, and the trick is almost always: \emph{unpack the definitions}.

\subsection*{Even and odd}

Let's nail down what ``even'' and ``odd'' actually mean, precisely enough to use in proofs.

\begin{definition}[Even and odd]
An integer $n$ is \emph{even}\index{even} if $n = 2k$ for some integer $k$. An integer $n$ is \emph{odd}\index{odd} if $n = 2k + 1$ for some integer $k$.
\end{definition}

These definitions are what we actually \emph{use} in proofs --- when we know something is even, we unpack the definition to get the $k$, and then we do algebra with $2k$. When we need to \emph{show} something is even, we exhibit a $k$ and show the equation holds. This is the ``there exists'' hidden inside these definitions: ``$n$ is even'' really means ``$\exists k \in \mathbb{Z}.\; n = 2k$.''

Let's do a worked example.

\textbf{Claim:} If $a$ and $b$ are both odd integers, then $a + b$ is even.

\textbf{Proof:} Let $a$ and $b$ be arbitrary odd integers. By the definition of odd, we have $a = 2j + 1$ for some integer $j$, and $b = 2k + 1$ for some integer $k$. Then
\[
  a + b = (2j + 1) + (2k + 1) = 2j + 2k + 2 = 2(j + k + 1).
\]
Since $j + k + 1$ is an integer, $a + b = 2m$ where $m = j + k + 1$, so $a + b$ is even by definition. \hfill $\square$

See the pattern? We assumed the hypothesis (both odd), unpacked the definitions (got $j$ and $k$), did some algebra, and then repacked the result into the definition of even. See the pattern? Assume the hypothesis, unpack the definitions, do the algebra, repack the result. That's the core of direct proof.

Let's do another one.

\textbf{Claim:} The sum of any even integer and any odd integer is odd.

\textbf{Proof:} Let $a$ be an arbitrary even integer and $b$ be an arbitrary odd integer. Then $a = 2j$ for some integer $j$ and $b = 2k + 1$ for some integer $k$. So
\[
  a + b = 2j + 2k + 1 = 2(j + k) + 1.
\]
Since $j + k$ is an integer, $a + b$ has the form $2m + 1$ where $m = j + k$, so $a + b$ is odd. \hfill $\square$

Same pattern, different details.

\subsection*{Divisibility}

Here's another definition that works exactly the same way.

\begin{definition}[Divisibility]
For integers $a$ and $d$ with $d \neq 0$, we say $d$ \emph{divides}\index{divisibility} $a$ (written $d \mid a$) if $a = dk$ for some integer $k$.
\end{definition}

So ``$3 \mid 12$'' because $12 = 3 \cdot 4$, and ``$5 \nmid 7$'' because there's no integer $k$ with $7 = 5k$.

Notice this has the same structure as evenness: it's an existential claim buried inside a definition. So proofs about divisibility work the same way --- unpack, do algebra, repack.

\textbf{Claim:} For all integers $a, b, c$: if $a \mid b$ and $b \mid c$, then $a \mid c$. (Transitivity of divisibility.)

\textbf{Proof:} Let $a, b, c$ be arbitrary integers with $a \mid b$ and $b \mid c$. By definition, $b = aj$ for some integer $j$, and $c = bk$ for some integer $k$. Then
\[
  c = bk = (aj)k = a(jk).
\]
Since $jk$ is an integer, we have $c = am$ where $m = jk$, so $a \mid c$. \hfill $\square$

\subsection*{Rational numbers}

One more definition in this style.

\begin{definition}[Rational number]
A real number $r$ is \emph{rational}\index{rational number} if $r = a/b$ where $a, b$ are integers and $b \neq 0$.
\end{definition}

Showing that a number is rational is really just an existence proof --- you need to exhibit the integers $a$ and $b$. For example, $0.75$ is rational because $0.75 = 3/4$, and we've exhibited $a = 3$, $b = 4$\footnote{Technically you could also use $a = 6$, $b = 8$, or any other representation. There's no uniqueness requirement in the definition, though you can get uniqueness if you insist on lowest terms.}.

\bigskip
\textbf{Exercise 1.} \label{ex:m5:1} Prove: the product of two rational numbers is rational. (Use the same unpack-algebra-repack pattern from the proofs above.)

\section{Counterexample}

Direct proofs let you \emph{prove} universal statements. But what if a universal statement is \emph{false}? How do you disprove it?

A universal statement ``for all $x$, $P(x)$'' is false if there exists \emph{even one} $x$ where $P(x)$ is false. That one example is called a \emph{counterexample}\index{counterexample}. This is just the negation of a universal: $\lnot(\forall x.\; P(x)) \equiv \exists x.\; \lnot P(x)$. We saw this equivalence back in Module 4!

So to disprove a universal, you just need to find one concrete value that breaks it. Let's see this in action.

\textbf{Claim to disprove:} For every integer $n$, if $n$ is odd then $(n^2 - 1)/2$ is odd.

\textbf{Disproof:} Counterexample: let $n = 1$. Then $n$ is odd, and $(1^2 - 1)/2 = 0/2 = 0$, which is even, not odd. So the claim is false. \hfill $\square$

That's the whole proof. One example. That's all you need.

Here's the important asymmetry: to \textbf{prove} a universal statement you need to work with an \emph{arbitrary} element --- you can't pick a specific one, because you need the argument to work for all of them. But to \textbf{disprove} a universal you just need \emph{one specific} counterexample. This asymmetry shows up everywhere and is worth internalizing.

\bigskip
\textbf{Exercise 2.} \label{ex:m5:2} Disprove: for every integer $n$, $n^2 > n$. (Find a counterexample and verify it.)

If you're a programmer, think of it this way: proving a function is correct means showing it works \emph{for all inputs}. Finding a bug means finding \emph{one bad input}. Testing is the art of finding counterexamples; formal verification is the art of proving universals. That's why testing can show the presence of bugs but never their absence\footnote{This is usually attributed to Dijkstra, and it's one of those quotes that sounds deep until you realize it's just the logical asymmetry between $\forall$ and $\exists$ wearing a trench coat.}.

\section{Existence proofs}

Counterexamples are about \emph{disproving} universals, but what about \emph{proving} existentials? If you want to show $\exists x : A.\; P(x)$, how do you actually do it?

The most straightforward way is a \emph{constructive existence proof}\index{existence proof}\index{constructive proof}: you produce a concrete witness and verify that it satisfies the predicate. This is exactly $\exists$-introduction from Module 4! You hand over a specific element $c : A$ along with a proof that $P(c)$ holds, and you're done.

\textbf{Claim:} There exists a real number $x > 1$ such that $2^x > x^{10}$.

\textbf{Proof:} Let $x = 100$. Then $x > 1$ and we need to check that $2^{100} > 100^{10}$. Well, $100^{10} = (10^2)^{10} = 10^{20}$. And $2^{10} = 1024 > 10^3$, so $2^{100} = (2^{10})^{10} > (10^3)^{10} = 10^{30}$. Since $10^{30} > 10^{20}$, we have $2^{100} > 100^{10}$. \hfill $\square$

That's it! We picked a value, plugged it in, and checked. The hard part isn't the proof itself --- it's figuring out \emph{which} value to try. But once you've found one that works, the proof is just computation.

A couple of things to note about existence proofs:
\begin{itemize}
  \item You don't need to find the \emph{best} or \emph{smallest} witness. Any witness that works is fine. We used $x = 100$ above, but $x = 59$ also works.\footnote{In fact the crossover point is somewhere around $x \approx 58.77$, but nobody asked us for the exact boundary. One working example is all we need.} The goal is existence, not optimization.
  \item The witness doesn't have to come from anywhere in particular. You can use trial and error, intuition, a calculator, whatever. As long as you \emph{verify} it rigorously in the proof, nobody cares how you found it.
  \item This is a \emph{constructive} proof because we actually built the thing we claimed exists. There are also \emph{non-constructive} existence proofs that show something must exist without ever telling you what it is---proof by contradiction can sometimes do this---but constructive proofs are generally preferred when you can get them, especially in computer science where you usually want to actually \emph{compute} the thing.\footnote{This is the constructive mathematician in your instructor peeking through again. Hi.}
\end{itemize}

It's worth making the computational significance of this distinction more explicit. A constructive existence proof hands you a \emph{witness}---a concrete object satisfying the property---and in programming terms that means it's a function that actually returns a value. A non-constructive proof, by contrast, tells you something exists without giving you any way to compute it: you know it's out there, but you can't write a program to find it. This is a real practical distinction, not a philosophical one. If you prove ``there exists an optimal schedule'' constructively, you have an algorithm; if you prove it non-constructively, you have a theorem but no code. This is the core reason constructive mathematics appeals to computer scientists---every constructive proof is, in a precise sense, a program. (This idea can be made fully rigorous via the Curry--Howard correspondence, which we'll gesture at later in the course, but even informally it's a useful way to think about what your proofs are actually \emph{doing}.)

\bigskip

Now that we have the basic direct proof toolkit for number theory, let's practice these same techniques in a completely different mathematical setting. The point is to show that the proof patterns we just learned --- unpack definitions, do algebra, repack --- are \emph{general}. They don't just work for even/odd and divisibility; they work for any mathematical structure with precise definitions.

I'm going to talk about \emph{groups}.

\section{What is a group?}
A group is one of the simplest structures in abstract algebra\footnote{I sometimes think of this as ``My First Abstract Algebra''---it's the training-wheels version of a much larger zoo of algebraic structures.}---an incredibly useful but very simple concept in mathematics that shows up over and over again. It's also simple enough to start with if you've never seen any abstract algebra before.

Here's the definition:

\begin{definition}[Group]
A \emph{group}\index{group} is a set $G$, an element $e : G$\index{identity element}, and a binary operation $* : G \to G \to G$ that satisfies the following axioms:

(associativity) $\forall a,b,c : G.\; (a * b) * c = a * (b * c)$

(left-identity) $\forall g : G.\; e * g = g$

(right-identity) $\forall g : G.\; g * e = g$

(inverses)\index{inverse} $\forall g : G.\; \exists g^{-1} : G$ such that both
\begin{enumerate}
  \item (left-inverse) $g^{-1} * g = e$, and
  \item (right-inverse) $g * g^{-1} = e$.
\end{enumerate}
\end{definition}

So what is this saying? A group is a triple of data: a base set $G$, a binary operation $*$ on the set, and an element $e$ that acts as the identity for the operation. Associativity tells us that repeated applications of the operation don't depend on how we parenthesize them, the identity axioms guarantee that $e$ really acts as an identity, and the inverse axioms guarantee that every element has an inverse that can be used to cancel it out.

Let's start by proving that some things \textbf{are} or are \textbf{not} groups.

\section{The integers under addition form a group}
Proposal: the set of integers, $\mathbb{Z}$, is a group with addition and an identity element of 0.

Proof:
To prove this we need to show all of the axioms hold for the integers. We're allowed to assume the standard definitions of addition, negation, and 0 since we're proving that these form a group not defining them formally. If you were doing this \emph{inside} a logic you'd have to first use the tools of the logic to give definitions for $\mathbb{Z}$ and for all the operations on it. This is essentially what the Natural Number Game requires.

Proof of associativity:
Let $a,b,c : \mathbb{Z}$. Then $(a+b)+c = a+(b+c)$ by the associativity of integer addition.

Proof of left-identity:
Let $n : \mathbb{Z}$, an arbitrary element, we need to show that $0 + n = n$. But this follows by the definition of addition on $\mathbb{Z}$!

Proof of right-identity:
Let $n : \mathbb{Z}$, an arbitrary element, we need to show that $n + 0 = n$. But this follows by the definition of addition on $\mathbb{Z}$!

Now to show the inverses we need to show that for every element of $\mathbb{Z}$ we have a proposed inverse-under-addition for this element, and \emph{then} we prove that the element really is both the left and right inverse.

The obvious candidate for every $n : \mathbb{Z}$ would be $-n$!

So now the theorems we need to prove become

$\forall n.\; (-n) + n = 0$

$\forall n.\; n + (-n) = 0$

but this is just a funnily written subtraction! And we already have defined subtraction to work this way, so our sub-theorems work automatically.

Thus we've shown that $(\mathbb{Z},+,0)$ is a group. \hfill $\square$

\bigskip
\textbf{Exercise 3.} \label{ex:m5:3} Is $(\mathbb{Z}, -, 0)$ a group (the integers with subtraction as the binary operation and $0$ as the proposed identity)? Either verify all four axioms or find one that fails.

\bigskip
\textbf{Exercise 4.}\label{ex:m5:4} Find the error in the following ``proof'' that every integer is even:

``Let $n$ be an integer. Suppose $n = 2k$ for some integer $k$. Then $n = 2k$, which is even by definition. Therefore $n$ is even.''

What proof technique error does this commit?

\section{The integers under multiplication do \emph{not} form a group}
We can \emph{disprove} that the integers with multiplication as the binary operation and 1 as the identity are a group easily. We just need to show one of the theorems cannot hold.

Let's show that inverses don't exist!

Consider the integer $2$. For $(\mathbb{Z}, \cdot, 1)$ to be a group, there would need to be some integer $n$ such that $2 \cdot n = 1$. But for any integer $n$, the product $2 \cdot n$ is always even\footnote{If you want a bonus challenge, prove this property of multiplication with natural induction! It's not too hard but is a good exercise}, and $1$ is odd. So there is no integer that can act as the multiplicative inverse of $2$.

Since we've found even one element without an inverse, the inverse axiom fails and it's not a group! \hfill $\square$

\section{Proof of uniqueness of identity}
We can also prove properties about \emph{all} groups, not just a specific group. Like let's prove the uniqueness of identity elements. In other words, that if there are two elements in a group that satisfy the identity axioms then they must be \textbf{equal}.

Theorem: Let $(G,*,e)$ be a group and let $e' : G$ satisfy the left and right identity properties. Then $e = e'$.

Proof:
We'll prove this by calculation with the properties of identity.

\begin{align*}
e &= e * e' && \text{(since $e'$ is a right-identity: $e * e' = e$)} \\
e * e' &= e' && \text{(since $e$ is a left-identity: $e * e' = e'$)}
\end{align*}
and thus $e = e'$. \hfill $\square$

\bigskip

This result tells us something important: the identity element of a group isn't just convenient notation---it's \emph{forced}. Any element that behaves like an identity must be the same as the one we started with. This kind of uniqueness result is typical in algebra: the axioms pin things down more tightly than you might expect.

\end{document}
